% !TEX root = ../main.tex
\chapter{Beam-Angle Calibration Analysis}

This chapter analyzes the combined forward and reverse calibration data stored in
\path{calibration_logs/combined_20260407_forward_reverse/combined_cleaned_calibration.csv}.
Those data are used here as the direct source for beam-angle calibration figures and
hysteresis statistics.

\section{Scope of the Calibration Data}

The combined CSV merges a forward sweep and an interpolated reverse sweep onto a common
AS5600 angle grid. Its role in this report is to show the measured static beam-angle
relationship and the amount of direction-dependent hysteresis. It is not the runtime
controller source-of-truth by itself.

The implemented controller still uses the same-session affine calibration documented in
\path{modeling.md} and \path{src/main.cpp}:

\begin{equation}
  \theta_{\text{cal,deg}} = 0.07666806 \cdot \theta_{\text{AS5600,deg}} - 23.28443907
\end{equation}

That distinction matters because the combined forward/reverse dataset shows larger
direction-dependent offset than the same-session calibration chosen for closed-loop
control.

\section{Calibration Summary}

The combined CSV contains \CalibrationPointCount{} grid points. The observed reverse
minus forward hysteresis is summarized by:

\begin{itemize}
  \item mean hysteresis: \CalibrationMeanHysteresisDeg{} degrees,
  \item RMS hysteresis: \CalibrationRmsHysteresisDeg{} degrees, and
  \item max hysteresis: \CalibrationMaxHysteresisDeg{} degrees.
\end{itemize}

\ReportTableFromFile
  {Summary statistics generated directly from the combined calibration CSV.}
  {tables/calibration_summary.tex}
  {tab:calibration_summary}

\section{Forward/Reverse Curve}

\ReportFigure[width=0.95\linewidth]
  {figures/calibration_curve.pdf}
  {Forward and reverse beam-angle calibration curves generated directly from the
  combined calibration CSV. The overlaid affine line is the runtime mapping used in the
  controller.}
  {fig:calibration_curve}

Figure~\ref{fig:calibration_curve} shows that the relationship is approximately affine
over the operating window but not perfectly single-valued when travel direction is
considered. The combined dataset therefore serves best as an offline characterization
tool.

\section{Hysteresis Plot}

\ReportFigure[width=0.95\linewidth]
  {figures/calibration_hysteresis.pdf}
  {Direction-dependent hysteresis, plotted as reverse minus forward beam angle versus
  unwrapped AS5600 angle.}
  {fig:calibration_hysteresis}

The hysteresis remains positive across the combined dataset, with a magnitude large
enough to justify separating the thesis-facing static characterization from the simpler
runtime affine map used for controller execution.
